%% =====================================================================
%%  T-14-01  The Automatic Response
%%  -------------------------------------------------------------------
%%  One idea per file. Core is mandatory. Slots in canonical order.
%%  Max 700 words. No layout commands. No filler. New source material
%%  for this entry goes in sources/<BOOK>/T-14-01.tex -- never by
%%  rewriting this file (docs/RULES.md R0.1).
%% =====================================================================
%% @kind: topic
%% @id: T-14-01
%% @title: The Automatic Response
%% @subtitle: A single cue can fire a whole routine
%% @chapter: c14
%% @order: 10
%% @status: stable
%% @sources: B4
%% @updated: 2026-09-19
%% @allow-rewrite: slot migration to the seven-question format (docs/RULES.md section 2)

\topic{T-14-01}{The Automatic Response}{A single cue can fire a whole routine}


\begin{core}
Most compliance is not decided. A salient feature of the situation --- a title, a queue, a deadline, a small gift --- fires a rehearsed response that runs without the substance being evaluated. Whoever produces the feature steers the response and never has to win an argument.
\end{core}

\begin{mechanism}
The routine exists because evaluating everything is unaffordable. Judging a request on its merits takes attention, information and time; most requests are ordinary, so a shortcut is normally correct. The shortcut is to find one feature that reliably predicts the whole and respond to the feature rather than to the thing.

That is efficient, and it is precisely what makes it exploitable: the feature can be manufactured far more cheaply than the condition it stands for. A credential is faster to fake than to earn, a crowd faster to assemble than to convince, a deadline faster to announce than to meet. Every principle in the rest of this book is a case of a predictable feature being produced on purpose, and the target's error is never stupidity --- it is that the shortcut was correct often enough to be worth keeping.
\end{mechanism}

\begin{conditions}
The same words on letterhead are read differently from the same words on a napkin, and the difference survives knowing perfectly well that letterhead costs nothing to print.
\end{conditions}

\begin{application}
  \item Identify the one feature you intend the other party to respond to, and make sure it is present before you make the request.
  \item Ask what it would cost you to fake that feature. If the answer is little, expect the other party to have faked it.
  \item Make your own cues expensive to counterfeit: a verifiable record, a named reference, a real queue. A cue that cannot be borrowed cannot be stolen.
\end{application}

\begin{feedback}
  \item You notice yourself agreeing before you have finished hearing the request.
  \item Your reason for saying yes is a single feature --- the letterhead, the logo, the queue --- rather than a weighing.
  \item The feature you relied on is cheap to produce and expensive to verify.
  \item Asked afterwards to reconstruct the decision, you cannot.
\end{feedback}

\begin{failure}
The shortcut is not a defect and cannot be removed without cost. A person who evaluates every cue from first principles cannot hire anyone, buy anything or cross a street, and is still wrong about most things --- only slower and more exhausting about it. The usable rule is asymmetric: keep the shortcuts where being wrong is cheap, and treat cheap-to-fake cues as worthless exactly where it is not. The entry's epilogue case is that the volume of cues now arriving per hour has risen while the time to evaluate each has fallen, which moves more decisions into the first category whether or not anyone intends it.
\end{failure}

\begin{countermeasures}
  \item Separate the cue from the claim, and ask the question the cue was answering in terms the cue cannot answer.
  \item Delay. Automatic routines depend on momentum, and momentum does not survive an hour.
  \item Price the cue. Anything cheap to produce and hard to verify carries no information and should be discounted to zero.
  \item Name the routine to yourself while it is running. Recognising a shortcut in use is most of disabling it.
\end{countermeasures}

\sources{B4}
\seesources
\seealso{T-14-02, T-01-01, T-18-01, T-23-01}
